% !TeX root = main.tex

\section{Inleiding}
In de hedendaagse topsport wordt het verschil tussen winst en verlies vaak bepaald door fracties van seconden en enkele centimeters. Waar tijdwaarneming zich vroeger beperkte tot het vastleggen van de finish, vereisen moderne analyses en live tracking een continue en accurate positiebepaling van atleten en voertuigen. \\

MYLAPS Sports Timing speelt hierin wereldwijd een belangrijke rol. Het van oorsprong Nederlandse bedrijf levert al decennialang de hardware en software voor duizenden evenementen, variërend van lokale hardloopwedstrijden tot de Olympische Spelen en grote motorsportkampioenschappen. De kern van MYLAPS breidt zich echter steeds verder uit. Naast het accuraat registreren van tijden, ontstaat er vanuit de markt een sterke behoefte aan robuuste, continue real-time analyse van deelnemers gedurende een gehele race. \\

In opdracht van MYLAPS wordt in dit afstudeerproject dan ook gewerkt aan de volgende generatie van deze volgsystemen, specifiek ontworpen voor de paardensport. Deze nieuwe systemen maken gebruik van Global Navigation Satellite Systems (GNSS) en kunnen in potentie zeer nauwkeurig zijn. MYLAPS maakt momenteel gebruik van Real-Time Kinematic (RTK) GNSS, waarbij de locatie tot op de centimeter nauwkeurig kan worden bepaald. In de dynamische praktijk van een renbaan ontbreekt echter vaak de garantie op een constant betrouwbaar signaal. Het risico dat een systeem onterecht een foutieve positie rapporteert, vormt een groot obstakel voor het gebruik in officiële wedstrijden. Daarom bestaat de wens om deze satellietdata te valideren en te corrigeren door gebruik te maken van extra sensoren, waarbij de Inertial Measurement Unit (IMU) centraal staat. Omdat deze sensor autonoom de oriëntatie en versnelling van het paard kan bijhouden, kan een veel robuustere en nauwkeurigere positiebepaling bereikt worden. \\

Dit Plan van Aanpak vormt de basis voor de uitvoering van het project en legt de grenzen vast waarbinnen het onderzoek en het ontwerp zullen plaatsvinden. Het doel van dit document is om de opdrachtomschrijving eenduidig te formuleren en een heldere leidraad te bieden voor het gehele traject. \\ 

In hoofdstuk 2 wordt de aanleiding van het project uitgelicht, waarna in hoofdstuk 3 en 4 meer context wordt gegeven over de gebruikte theorie. In hoofdstuk 5 worden de probleemstelling en doelstelling uiteengezet. Vervolgens worden in hoofdstuk 6 de leerdoelen uitgewerkt en in hoofdstuk 7 de technische specificaties en randvoorwaarden gedefinieerd. Dit vormt de basis voor het ontwerpaanpak in hoofdstuk 8. Het document sluit af met hoofdstuk 9, waarin het testplan wordt beschreven waarmee de werking van het eindresultaat wordt gevalideerd.



\newpage